\appendix
\part{Anhang}

\newcommand{\lit}[3]{\vspace{5pt}\noindent\textbf{[#1]} #2\par\nopagebreak\vspace{1pt}{\small\itshape\color{grau} #3}\par\vspace{2pt}}

\chapter{Formelsammlung und Übersicht}

\section{Die wichtigsten Formeln}

\begin{merksatz}
\boxtitel{Grundlagen}
\begin{center}
{\small
\renewcommand{\arraystretch}{1.55}
\begin{tabular}{@{}P{4.2cm}P{4.4cm}@{}}
\toprule
$\sum_{v} \deg(v) = 2m$ & Handschlaglemma (1.3)\\
Anzahl ungerader Grade ist gerade & (1.4)\\
$m \le \binom{n}{2} = \frac{n(n-1)}{2}$ & höchstmögliche Kantenzahl\\
$\deg_{\overline G}(v) = n-1-\deg_G(v)$ & Komplement\\
$m \ge n - c$ & $c$ Komponenten (2.8, 4.3)\\
$(A^k)_{ij}$ & Kantenzüge der Länge $k$ (3.4)\\
$\frac16\operatorname{spur}(A^3)$ & Anzahl der Dreiecke (3.5)\\
\bottomrule
\end{tabular}}
\end{center}
\end{merksatz}

\begin{merksatz}
\boxtitel{Bäume und Spannbäume}
\begin{center}
{\small
\renewcommand{\arraystretch}{1.55}
\begin{tabular}{@{}P{4.2cm}P{4.4cm}@{}}
\toprule
Baum $\iff$ zusammenhängend, kreisfrei & fünf Fassungen (4.2)\\
$m = n-1$ & Baum\\
$m = n - c$ & Wald mit $c$ Komponenten\\
$\ge 2$ Blätter für $n \ge 2$ & (4.3)\\
$n^{\,n-2}$ & Bäume auf $n$ Ecken (5.5)\\
$\frac{(n-2)!}{\prod_i (d_i-1)!}$ & mit vorgeschriebenen Graden (5.6)\\
$\det L_i$ & Spannbäume (5.8)\\
$r^{s-1}s^{r-1}$ & Spannbäume von $K_{r,s}$\\
\bottomrule
\end{tabular}}
\end{center}
\end{merksatz}

\begin{merksatz}
\boxtitel{Durchlaufen, Zeichnen, Färben}
\begin{center}
{\small
\renewcommand{\arraystretch}{1.55}
\begin{tabular}{@{}P{4.2cm}P{4.4cm}@{}}
\toprule
Eulerkreis $\iff$ alle Grade gerade & zusammenhängend (6.4)\\
Eulerzug $\iff$ genau zwei ungerade & (6.5)\\
$\deg(u)+\deg(v)\ge n$ $\Rightarrow$ hamiltonsch & Ore (7.3)\\
$\deg(v) \ge n/2$ $\Rightarrow$ hamiltonsch & Dirac (7.4)\\
$c(G-S) \le |S|$ & notwendig (7.2)\\
$n - m + f = 2$ & Polyederformel (8.3)\\
$m \le \frac{g(n-2)}{g-2}$ & Taille $g$, planar (8.4)\\
$\delta(G) \le 5$ & planar (8.6)\\
$\chi \ge \omega$, $\chi \ge n/\alpha$ & (9.3)\\
$\chi \le \Delta + 1$ & gierig (9.4)\\
$P(G,k) = P(G-e,k) - P(G/e,k)$ & (10.3)\\
$P(C_n,k) = (k-1)^n + (-1)^n(k-1)$ & (10.6)\\
\bottomrule
\end{tabular}}
\end{center}
\end{merksatz}

\begin{merksatz}
\boxtitel{Paarungen, Flüsse, Extremales, Spektren}
\begin{center}
{\small
\renewcommand{\arraystretch}{1.55}
\begin{tabular}{@{}P{4.2cm}P{4.4cm}@{}}
\toprule
$|N(S)| \ge |S|$ & Hall (11.4)\\
größte Paarung $=$ kleinste Überdeckung & Kőnig, bipartit (11.7)\\
größter Fluss $=$ kleinster Schnitt & (12.5)\\
Wege $=$ trennende Menge & Menger (12.8)\\
$m \le n^2/4$ & dreiecksfrei, Mantel (13.1)\\
$m \le (1-\frac1r)\frac{n^2}{2}$ & ohne $K_{r+1}$, Turán (13.3)\\
$R(3,3) = 6$ & (13.4)\\
$R(s,s) > 2^{s/2}$ & probabilistisch\\
$\frac{2m}{n} \le \lambda_1 \le \Delta$ & (14.3)\\
$x^{\!\top}Lx = \sum_{uv \in E}(x_u-x_v)^2$ & (14.5)\\
Vielfachheit von $0$ $=$ $c$ & Laplace (14.6)\\
$p \approx 1/n$ & Schwelle Dreieck (14.8)\\
$p \approx \ln n / n$ & Schwelle Zusammenhang (14.9)\\
\bottomrule
\end{tabular}}
\end{center}
\end{merksatz}

\section{Die Landkarte des Bandes}

\begin{merksatz}
\boxtitel{Welcher Satz wo steht --- und worauf er beruht}
\begin{center}
{\small
\renewcommand{\arraystretch}{1.42}
\begin{tabular}{@{}P{3.4cm}P{1.4cm}P{3.4cm}@{}}
\toprule
\textbf{Satz} & \textbf{Kap.} & \textbf{Beruht auf}\\
\midrule
Handschlaglemma & $1$ & Doppelzählung\\
Kantenzug $\to$ Weg & $2$ & Extremalprinzip\\
Bipartit $\iff$ kein ungerader Kreis & $2$ & Abstandsfärbung\\
Brücke $\iff$ kein Kreis & $2$ & ---\\
Potenzen zählen Kantenzüge & $3$ & Induktion\\
Fünf Baumfassungen & $4$ & Ringschluss, Satz 2.14\\
Kruskal & $5$ & Austauschargument\\
Cayley & $5$ & Prüfer-Bijektion\\
Matrix-Baum-Satz & $5$ & Cauchy--Binet\\
Euler/Hierholzer & $6$ & Verschmelzen\\
Ore & $7$ & maximales Gegenbeispiel\\
Petersen nicht hamiltonsch & $7$ & drei Fälle\\
Polyederformel & $8$ & Induktion, Jordan\\
Kantenschranke & $8$ & Doppelzählung\\
Fünffarbensatz & $9$ & Satz 8.6, Kempe-Kette\\
Chromatisches Polynom & $10$ & Mengenaufteilung\\
Hall & $11$ & Induktion, zwei Fälle\\
Kőnig & $11$ & alternierende Wege\\
Max-Fluss-Min-Schnitt & $12$ & Restnetzwerk\\
Menger & $12$ & Übersetzung\\
Mantel & $13$ & Quadratsummen\\
Turán & $13$ & Induktion in $r$-Schritten\\
$R(3,3)=6$ & $13$ & Schubfach\\
Laplace-Nullstellen & $14$ & Satz 14.5\\
\bottomrule
\end{tabular}}
\end{center}
\end{merksatz}

\begin{merksatz}
\boxtitel{Was in diesem Band \kursiv{nicht} bewiesen wurde}
Vier Ergebnisse wurden benutzt, ohne sie herzuleiten. Alle vier sind an ihrer Stelle ausdrücklich als solche gekennzeichnet.

\begin{center}
{\small
\renewcommand{\arraystretch}{1.5}
\begin{tabular}{@{}P{3.0cm}P{1.4cm}P{3.8cm}@{}}
\toprule
\textbf{Ergebnis} & \textbf{Kap.} & \textbf{Wo es steht}\\
\midrule
Jordanscher Kurvensatz & $8$ & Band 8 dieser Reihe\\
Satz von Kuratowski & $8$ & Diestel\zit{7}, Kap. 4\\
Satz von Brooks & $9$ & Diestel\zit{7}, Kap. 5\\
Satz von Tutte & $11$ & Diestel\zit{7}, Kap. 2\\
Spektralsatz & $14$ & Band 2, Teil V\\
Cauchy--Binet & $5$ & Band 2, Teil IV\\
\bottomrule
\end{tabular}}
\end{center}

\begin{center}
\begin{minipage}{0.88\textwidth}
\kursiv{Sechs Auslassungen auf zweihundertfünfzig Seiten. Jede ist im Text benannt, mit Begründung, warum sie dort steht. Das ist die einzige Form, in der die Regel dieses Werkes --- nichts ist selbstverständlich --- gebrochen werden darf: offen und mit Quellenangabe.}
\end{minipage}
\end{center}
\end{merksatz}

\section{Zeittafel}

\begin{merksatz}
\boxtitel{Von Königsberg bis heute}
\begin{center}
{\small
\renewcommand{\arraystretch}{1.38}
\begin{tabular}{@{}P{1.4cm}P{2.8cm}P{4.0cm}@{}}
\toprule
\textbf{Jahr} & \textbf{Person} & \textbf{Ereignis}\\
\midrule
$1736$ & Euler & Königsberger Brücken\\
$1750$ & Euler & Polyederformel\\
$1847$ & Kirchhoff & Inzidenzmatrix, Spannbäume\\
$1852$ & Guthrie & Vierfarbenfrage\\
$1857$ & Cayley & Begriff des Baumes\\
$1859$ & Hamilton & das Icosian-Spiel\\
$1873$ & Hierholzer & Beweis des Eulersatzes\\
$1875$ & Cayley & Zählung der Alkane\\
$1878$ & Sylvester & das Wort \enquote{Graph}\\
$1879$ & Kempe & fehlerhafter Vierfarbenbeweis\\
$1889$ & Cayley & $n^{n-2}$\\
$1890$ & Heawood & Fünffarbensatz\\
$1891$ & Petersen & der Petersen-Graph\\
$1907$ & Mantel & dreiecksfreie Schranke\\
$1912$ & Birkhoff & chromatisches Polynom\\
$1916$ & Kőnig & bipartit, Paarungen\\
$1918$ & Prüfer & der Prüfer-Code\\
$1926$ & Borůvka & minimaler Spannbaum\\
$1927$ & Menger & Wege und Trennmengen\\
$1930$ & Kuratowski & Planaritätskriterium\\
$1930$ & Ramsey & Ramseytheorie\\
$1931$ & Kőnig, Egerváry & Paarung $=$ Überdeckung\\
$1935$ & Hall & Repräsentantensysteme\\
$1936$ & Kőnig & erstes Lehrbuch\\
$1937$ & Wagner & Minorenfassung\\
$1941$ & Turán & extremale Schranke\\
$1943$ & Hadwiger & die Vermutung\\
$1947$ & Erdős & probabilistische Methode\\
$1947$ & Tutte & perfekte Paarungen\\
$1952$ & Dirac & Hamiltonbedingung\\
$1956$ & Ford, Fulkerson & Flusssatz\\
$1956$ & Kruskal & Spannbaumverfahren\\
$1957$ & Berge & vergrößernde Wege\\
$1959$ & Erdős, Rényi & Zufallsgraphen\\
$1960$ & Ore & Verallgemeinerung von Dirac\\
$1965$ & Edmonds & Paarungen allgemein\\
$1972$ & Karp & NP-Vollständigkeit\\
$1975$ & Szemerédi & Regularitätslemma\\
$1976$ & Appel, Haken & Vierfarbensatz\\
$1988$ & Lubotzky u.\,a. & Ramanujan-Graphen\\
$2004$ & Robertson, Seymour & Graphenminoren-Satz\\
$2005$ & Gonthier & Vierfarbensatz formalisiert\\
$2012$ & Huh & Log-Konkavität\\
$2016$ & Babai & Isomorphie quasipolynomiell\\
$2022$ & Chen u.\,a. & Flüsse fast linear\\
$2023$ & Campos u.\,a. & Ramsey-Schranke verbessert\\
$2026$ & Diskin u.\,a. & Perkolationsübergang\\
\bottomrule
\end{tabular}}
\end{center}
\end{merksatz}

\section{Glossar}

\begin{merksatz}
\boxtitel{Begriffe von A bis Z}
\begin{center}
{\small
\renewcommand{\arraystretch}{1.32}
\begin{longtable}{@{}P{3.0cm}P{5.6cm}@{}}
\toprule
\textbf{Begriff} & \textbf{Bedeutung}\\
\midrule
\endhead
Adjazenzmatrix & $n\times n$-Matrix mit Einsen an den Kantenstellen (3.2)\\
Baum & zusammenhängend und kreisfrei (4.1)\\
Bipartit & Ecken in zwei Gruppen ohne innere Kanten (2.11)\\
Blatt & Ecke vom Grad $1$\\
Brücke & Kante, deren Entfernung trennt (2.13)\\
Chromatische Zahl & kleinste Farbenzahl $\chi$ (9.1)\\
Clique & paarweise benachbarte Ecken (9.2)\\
Durchmesser & größter vorkommender Abstand (2.9)\\
Ecke & Element von $V$ (1.1)\\
Eulerkreis & benutzt jede Kante genau einmal (6.2)\\
Expander & dünn und dennoch fest verbunden (14)\\
Gerichteter Graph & Kanten mit Richtung (12.1)\\
Grad & Anzahl der Nachbarn (1.2)\\
Gradfolge & sortierte Liste aller Grade (1.5)\\
Graphon & Grenzwert einer Graphenfolge (15)\\
Hamiltonkreis & besucht jede Ecke genau einmal (7.1)\\
Induzierter Teilgraph & alle Kanten zwischen gewählten Ecken (1.6)\\
Invariante & bei Isomorphie gleichbleibende Größe (1.8)\\
Isomorph & strukturgleich (1.7)\\
Kante & zweielementige Eckenmenge (1.1)\\
Kantenzug & Folge benachbarter Ecken (2.1)\\
Kempe-Kette & Komponente aus zwei Farbklassen (9)\\
Komponente & Äquivalenzklasse der Erreichbarkeit (2.7)\\
Kreis & geschlossener Weg der Länge $\ge 3$ (2.2)\\
Laplace-Matrix & $L = D - A$ (5.7)\\
Minor & durch Löschen und Zusammenziehen (8)\\
Multigraph & mit Schlingen und Parallelen (6.1)\\
Paarung & Kanten ohne gemeinsame Ecke (11.1)\\
Perfekte Paarung & sättigt jede Ecke (11.1)\\
Planar & kreuzungsfrei zeichenbar (8.1)\\
Prüfer-Code & Folge zu einem beschrifteten Baum (5)\\
Schnitt & Aufteilung mit $s$ links, $t$ rechts (12.3)\\
Schnittecke & Ecke, deren Entfernung trennt (2.13)\\
Spannbaum & Baum auf allen Ecken (5.1)\\
Spektrum & Liste der Eigenwerte (14.2)\\
Taille & Länge des kürzesten Kreises (2.9)\\
Turán-Graph & vollständig mehrteilig (13.2)\\
Unabhängige Menge & Ecken ohne Kante dazwischen (9.2)\\
Vergrößernder Weg & alternierend, freie Enden (11.2)\\
Wald & kreisfrei, evtl. mehrere Bäume (4.1)\\
Weg & Kantenzug ohne Eckenwiederholung (2.2)\\
Wurzelbaum & Baum mit ausgezeichneter Ecke (4.6)\\
Zufallsgraph & $G(n,p)$ (14.7)\\
\bottomrule
\end{longtable}}
\end{center}
\end{merksatz}
